% !TEX root = ../../../main.tex
% Appendix entry: interrogation.cycle_shapes -- "Cycle shapes" (Appendix A).
\entryfield{Algorithm}
\code{_run} in \code{Adapters/interrogation_cycle_shapes.py} checks the optional
\emph{waveform} side input (a \texttt{Signal} of the chain's length and sampling rate) and
calls \code{cycle_shapes} in \code{Working/interrogation/cycle_shapes.py}. The method is the
back half of \citet{colominas2023iterative}'s iterative warping and clustering, written from
the paper (the authors' MATLAB carries no licence and was not read). Every adaptation is
argued, with measurements, in that module's docstring; the steps:

\begin{enumerate}
\item \emph{Event times and sequences.} An event's anchor is its \code{onset_idx} (with
  \emph{anchor} = \texttt{onset} and the column present) or its span start. Sequences come
  from \code{train_id} if present (events in no train open no cycle,
  \code{in_no_train}), else from \code{group}, else the whole span (\code{sequences_of}).
  An event with a NaN anchor is \code{no_anchor}. Within a sequence, sorted by anchor,
  cycle $k$ runs from event $k$'s anchor to event $k{+}1$'s; the last event opens none.
\item \emph{Gaps and splits.} With $\tilde\ell_k$ the median of the 11 cycle lengths around
  cycle $k$: a cycle longer than \emph{max\_cycle\_s}, or (at 0) than $1.6\,\tilde\ell_k$,
  is \code{gap_long}; one shorter than \emph{min\_cycle\_s}, or (at 0) than $0.4\,\tilde\ell_k$,
  or than 8 samples, is \code{too_short}. Both are left out and counted; over 10\,\% gaps in
  a sequence raises a warning.
\item \emph{Cut} (\code{prepare_cycle}) from the \emph{waveform} if bound, else from the
  chain's signal. More than \emph{max\_nan\_frac} NaN: \code{missing_values}; fewer are
  interpolated. With \emph{cycle\_detrend} = \texttt{chord}, subtract the straight line
  between the pre-drop levels at the two ends, each the median of the 2\,\% of samples (at
  least 2) just before that anchor. Zero peak-to-peak: \code{flat}. \code{cycle_amplitude} is
  the peak-to-peak of the detrended cut.
\item \emph{Warp}: resample to \emph{n\_points} (256) and z-normalise.
\item \emph{Refine} (\code{refine}): align every cycle to the pointwise median of the stack by
  zero-padded cross-correlation, the shift bounded by \emph{max\_shift\_frac} of the cycle;
  move the anchors by that shift (never more than that fraction from the detector's
  anchor), re-cut and re-warp. A round is kept only while the SVD entropy (eq.~17 of the
  paper) falls by more than its own bootstrap SD over 5 resamples; at most \emph{max\_iter}
  rounds.
\item \emph{Align} (\code{align_stack}): bounded cross-correlation to the median template,
  then to the stack's medoid (the cycle with the highest summed correlation to the others);
  vacated samples take the row's edge value; re-z-normalise. Not the paper's cyclic shifts:
  a fungal cycle does not wrap round.
\item \emph{Choose groups} (\code{choose_groups}) on the refined, warped (not aligned) stack.
  $K = 1$ unless the best two-way split (2-means, 50 starts, on the top 32 principal
  components) passes all of:
  \begin{itemize}
  \item its cluster index (within over total sum of squares) is below that of the same split
    on every null stack and more than 3 robust SDs below their median; the null stacks are
    the background pseudo-cycles (step 9), and when there is a background, also per-cycle
    phase-randomised copies of the real cycles (both must be beaten); when the background
    fell back, the phase copies alone;
  \item the cycles projected onto the line through the two centroids are two modes: they
    split better than uniform samples of the same size (\code{bimodality_test});
  \item each half has at least \emph{min\_cycles} cycles.
  \end{itemize}
  Each group is split again only while its own split passes the same tests at a level
  scaled by its size \citep{kimes2017statistical}, up to \emph{k\_max} groups, settled by one
  Lloyd pass. Groups are numbered 1..$K$ by size. SigClust \citep{liu2008statistical},
  without its background-noise floor, and the Calinski--Harabasz choice are reported but do
  not decide.
\item \emph{Change points and order}: a used cycle whose group differs from the previous used
  cycle's is a change (\code{cycle_change}). The order test counts runs of equal groups
  against \emph{n\_perm} shuffles of the cycle order: $p = (1 + \#\{\text{shuffles with as few
  runs}\})/(1 + n_{\text{perm}})$. It tests order, not shape.
\item \emph{Coherence and its verdict.} SVD entropy (\code{svd_entropy}, eq.~17, and the
  $\sigma^2$ variant, shown first) of the whole stack and of each group, at three points:
  fixed-length cuts, after warping, after alignment. The verdict null is \emph{null\_draws}
  draws of pseudo-cycles cut where no event lies, with lengths drawn from the real cycles,
  through the same chord, warp, refinement and alignment; a group is judged against the null
  group of its own entropy rank from the same $K$-means on each draw. A value beats the null
  when it is below every draw and more than 3 robust SDs (1.4826 MAD, floored at 0.002 of
  the median) below the draws' median. A group's key is unmarked only when eq.~17 and
  $\sigma^2$ both beat it; otherwise \texttt{\textasciitilde noise\_like}. If more than half
  the pseudo-cycles found no free stretch (a chain that fills its span), or the whole stack
  beats the background but not the phase copies of its own cycles (a background of another
  colour), every group is \texttt{\textasciitilde untested}.
\item \emph{Plateau measures} (\code{plateau}), on each warped, chord-detrended cycle after a
  5-point running median: \code{cycle_floor_frac} is the share of the cycle from the first
  sample in the lowest \emph{floor\_band} of its range to the last one before the climb's
  top; \code{cycle_half_recovery_frac} is where (0--1) it first gets half way back from its
  deepest point to the pre-drop level, NaN if never.
\end{enumerate}

\entryfield{Assumptions}
\begin{itemize}
\item \emph{Units.} Not unit-aware: \code{cycle_amplitude} is in the stored unit of the trace
  the cycles were cut from. Everything else is z-normalised or a fraction. Lengths and
  shifts are in seconds via $f_s$.
\item \emph{Signal shape.} Built for chains of drops with a slow climb between them. A
  detector usually runs on a short-detrended trace that flattens the climb (from a 600\,s
  rolling-mean trace, synthetic sharkfins read half back at 0.16 of the cycle against 0.69
  from the raw trace); bind \emph{waveform} to the root signal. An unbound waveform is
  warned.
\item \emph{Missing values}: per cycle, as in step 3. The background windows avoid NaN.
\item \emph{Edges.} Cycles exist only between events, so the stretch before the first event
  and after the last is never a cycle; it can be background.
\item \emph{Count.} Under \emph{min\_cycles} usable cycles in a sequence: measures but no
  group (\code{too_few_cycles}); under twice that: $K = 1$ by construction.
\item \emph{Seeds.} One \emph{seed} drives every random draw (nulls, $k$-means starts, stack
  sample).
\end{itemize}

\entryfield{Outputs}
A \texttt{SpanSet} with the spans and upstream columns unchanged, plus, on the event that
opens each cycle: \code{cycle_index}, \code{cycle_length_s}, \code{cycle_amplitude},
\code{cycle_group} ($-1$ = none), \code{cycle_change}, \code{cycle_shift_s},
\code{cycle_corr}, \code{cycle_floor_frac}, \code{cycle_half_recovery_frac} and
\code{cycle_omit_reason} (\code{""} when used). Index, group, change, shift, correlation and
omit reason are bookkeeping columns. With \emph{write\_group}, \code{group} is the group key
(\code{t<train>:<g>} with several trains, \texttt{-1} for events whose cycle was left out or
that open none).

\code{meta} (made JSON-safe) carries per-sequence results (counts, omissions, iteration
history, $K$ and its tests, the order test, a Kendall trend of the floor fraction over time),
the coherence values and every null draw, the verdicts, \code{exemplar_events} (each group's
medoid cycle, a member), the \code{cycle_stack} drawing payload, \code{warnings},
\code{rules} and a summary line.

\entryfield{Cost}
Superlinear in the number of cycles: the adapter's cost model uses exponent 1.4, measured at
7.3\,s for 98 cycles and about 55\,s for 400 at the defaults on a loaded machine. Most of it is
the null: \emph{null\_draws} background stacks each refined and aligned, plus the $K$ nulls
and \emph{n\_perm} shuffles. The estimate assumes one cycle per 30 minutes of span. The run
can be cancelled between steps.

\entryfield{Failure modes}
\begin{itemize}
\item \emph{Long cycles lost.} The automatic gap rule removes the longest cycles first: true
  cycles lost at a jitter CV of 0.3 / 0.5 are about 5.5\,\% / 17.5\,\% (module docstring); on
  rec 3 it left out 21\,\% of cycles (\code{W3-evaluation.md}). A bias against long plateaus.
\item \emph{Low power for a minority shape.} A random walk at 30\,\% noise with 80/20 or
  90/10 shape proportions reads $K = 1$ (module docstring).
\item \emph{Untested, not noise-like.} A chain that fills its span, or sits on a background
  of a different colour from its quiet margins, gets \texttt{\textasciitilde untested}
  everywhere; on a drifting background at 30\,\% noise that check raised a false alarm on a
  real two-shape chain.
\item \emph{One drifting family.} A steadily changing shape stays one group by design; the
  drift shows only in the per-cycle measures and the floor-fraction trend.
\item \emph{The detrended trace.} With \emph{waveform} unbound on a short-detrended chain,
  floor and half-recovery fractions describe the detrended shape, not the recording's.
\end{itemize}

\entryfield{Validation}
\code{tests/test_cycle_shapes.py} (40 tests) checks answers on synthetic chains: two shapes
give two groups with one change at the switch and beat the background on both entropies; a
lengthening plateau stays one group with its drift in the floor fraction; chain-timed noise on
a random walk is one group in every seed; Poisson events on coloured noise give one
noise-like group; the same chain at 1 and 10\,Hz gives the same groups; missed and split
detections are caught; floor fractions land where expected; a span-filling chain is
untested, not noise-like. On real recordings (\code{docs/prompts/blocks/reports/W3-evaluation.md})
the \texttt{chain\_cycles} template beat its background null on all three long chains
(rec 3, rec 6, rec 542) and shipped; all three gave $K = 1$. The registry description still
says ``not yet run on a real recording''; that is out of date.

\entryfield{References}
\citet{colominas2023iterative} (the method's back half); \citet{liu2008statistical}
(SigClust, reported only); \citet{kimes2017statistical} (the sequential rule for splitting a
hierarchy). Sources: \code{docs/research/analysis-blocks-primary-sources.md}, section 2.

\entryfield{Tests}
\code{tests/test_cycle_shapes.py}, \code{tests/test_adapter_cycle_shapes.py} (columns,
\emph{write\_group}, suffixes, the side input checked and used, JSON meta, the estimate),
\code{tests/test_template_chain_cycles.py} (the template end to end, with and without the
waveform bound), \code{tests/test_webui_w3f_ui.py}, \code{tests/test_webui_run_routing.py}.
